% Metodo k.p con spin-orbita (§3): (a) livelli in Gamma, (b) bande vicino a Gamma (formule della slide 41)
% (b): parametri di tipo GaAs, Eg = 1.52 eV, Delta = 0.34 eV, E_P = 20 eV; ascissa k in 1/angstrom
\begin{tikzpicture}
% ---------------------------------------------------------------- (a) livelli
\begin{scope}[x=1cm, y=1.55cm]
 \def\lw{1.15}
 \node[font=\small, anchor=west] at (-0.25,2.6) {(a) livelli in $\Gamma$};
 % senza spin-orbita
 \draw[very thick, vA] (0,1.52) -- (\lw,1.52);
 \draw[very thick, vB] (0,-0.113) -- (\lw,-0.113);
 \node[font=\scriptsize, align=center, anchor=south] at (0.575,1.55) {$s$, $L=0$\\[-1pt]$\times2$};
 \node[font=\scriptsize, align=center, anchor=north] at (0.575,-0.16) {$p$, $L=1$\\[-1pt]$\times6$};
 \node[font=\scriptsize, black!60, anchor=north] at (0.575,-0.95) {senza $H_{\mathrm{so}}$};
 % con spin-orbita
 \draw[very thick, vA] (2.6,1.52) -- (2.6+\lw,1.52);
 \draw[very thick, vB] (2.6,0) -- (2.6+\lw,0);
 \draw[very thick, color=vC] (2.6,-0.34) -- (2.6+\lw,-0.34);
 \node[font=\scriptsize, anchor=west] at (2.65+\lw,1.52) {$J=\frac12$ ($\times2$)};
 \node[font=\scriptsize, anchor=west, align=left] at (2.65+\lw,0.09) {$J=\frac32$ ($\times4$): $hh$, $lh$};
 \node[font=\scriptsize, anchor=west] at (2.65+\lw,-0.4) {$J=\frac12$ ($\times2$): $so$};
 \node[font=\scriptsize, black!60, anchor=north] at (3.175,-0.95) {con $H_{\mathrm{so}}$};
 \draw[black!40, densely dashed] (\lw,1.52) -- (2.6,1.52);
 \draw[black!40, densely dashed] (\lw,-0.113) -- (2.6,0);
 \draw[black!40, densely dashed] (\lw,-0.113) -- (2.6,-0.34);
 \draw[<->, black!70] (2.42,0) -- (2.42,-0.34) node[midway, left, font=\scriptsize] {$\Delta$};
 \draw[<->, black!70] (1.72,0) -- (1.72,1.52) node[midway, fill=white, inner sep=1pt, font=\scriptsize] {$\Eg$};
 \draw[black!25, densely dotted] (1.72,0) -- (2.6,0);
\end{scope}
% ---------------------------------------------------------------- (b) bande
\begin{axis}[at={(7.4cm,-2.15cm)}, anchor=south west, width=7.0cm, height=7.0cm,
  xmin=-0.13, xmax=0.13, ymin=-1.0, ymax=2.45, samples=101, domain=-0.13:0.13,
  axis lines=left, xtick={0}, xticklabels={$\Gamma$}, ytick={-0.34,0,1.52}, yticklabels={$-\Delta$,$0$,$\Eg$},
  y tick label style={font=\scriptsize}, clip=false, enlargelimits=false,
  xlabel={$k$}, ylabel={$E$}, every axis x label/.style={at={(axis description cs:1,0)}, anchor=west, xshift=3pt},
  every axis y label/.style={at={(ticklabel* cs:1.0)}, anchor=south, rotate=0}]
 \draw[black!25, densely dashed] (axis cs:0,-1.0) -- (axis cs:0,2.45);
 \draw[black!20] (axis cs:-0.13,0) -- (axis cs:0.13,0);
 % hbar^2/2m0 = 3.81 eV A^2 ; P^2 = E_P hbar^2/2m0 = 20*3.81 eV^2 A^2
 \addplot[very thick, vA, domain=-0.125:0.125] {1.52 + 3.81*x^2 + 20*3.81*x^2/3*(2/1.52 + 1/1.86)};
 \addplot[very thick, vB] {-3.81*x^2};
 \addplot[very thick, vB, densely dashed, domain=-0.13:0.13] {-3.81*x^2 - 2*20*3.81*x^2/(3*1.52)};
 \addplot[very thick, color=vC, domain=-0.13:0.13] {-0.34 - 3.81*x^2 - 20*3.81*x^2/(3*1.86)};
 \node[font=\scriptsize, vA, anchor=west] at (axis cs:0.0,2.38) {$c$ ($s$, $J=\frac12$)};
 \node[font=\scriptsize, vB, anchor=south] at (axis cs:0.105,0.0) {$hh$};
 \node[font=\scriptsize, vB, anchor=east] at (axis cs:-0.105,-0.36) {$lh$};
 \node[font=\scriptsize, color=vC, anchor=north] at (axis cs:0.085,-0.5) {$so$};
 \node[font=\small, anchor=south west] at (axis cs:-0.085,2.5) {(b) bande vicino a $\Gamma$};
\end{axis}
\end{tikzpicture}
